\section{Capacity, from the measured constants only}
\label{sec:capacity}

The constants in \S\ref{sec:state}--\S\ref{sec:isolation} determine a capacity
model. Everything in this section is \modeled{} --- arithmetic over measured
inputs and stated assumptions --- and nothing in it was observed at the scales
named. We give it because the arithmetic is the point of the decomposition, and
we mark it so that no reader mistakes it for a result.

\subsection{Storage is set by population, and is negligible}

Dormant agents are \rFleetPer{}\,bytes each, so storage is exactly linear in how
many agents exist:

\begin{center}
\begin{tabular}{lrr}
\toprule
\textbf{Registered agents} & \textbf{On disk} & \textbf{At list price/mo} \\
\midrule
$10^6$  & 454.6\,MiB & \$\rFleetStorage{} \meas{}/\modeled{} \\
$10^8$  & 44.4\,GiB  & \$0.89 \modeled{} \\
$10^9$  & 444.2\,GiB & \$8.90 \modeled{} \\
\bottomrule
\end{tabular}
\end{center}

A billion registered agents is under half a terabyte and under ten dollars a
month of object storage. Storage is not the constraint and it will not become
the constraint; \S\ref{sec:limits} names the thing that does.

\subsection{Memory is set by concurrency, and the working set dominates}

Live agents cost memory. The measured floor is \rGoPer{}\,bytes, but a working
agent holds context, so the honest model is parameterized by that working set
rather than by the floor. On a 128\,GiB node:

\begin{center}
\begin{tabular}{lrrr}
\toprule
\textbf{Per-agent working set} & \textbf{Live agents/node} & \textbf{$10^9$ reg.\ @ 1\%} & \textbf{@ 0.1\%} \\
\midrule
\rGoPer{}\,B (measured floor)  & 228\,M & 1 node & 1 node \\
10\,KiB                        & 13.4\,M & 1 node & 1 node \\
64\,KiB                        & 2.1\,M  & 5 nodes & 1 node \\
1\,MiB                         & 131\,k  & 77 nodes & 8 nodes \\
\bottomrule
\end{tabular}
\end{center}

\modeled{} throughout. The left column is the one that matters and it is the one
this paper did not measure: a fleet's cost is set by how much context its agents
hold while running, not by the residency floor of the primitive holding them.
What the measurement establishes is that the \emph{floor} contributes nothing ---
at \rGoPer{}\,bytes it is invisible against any realistic working set --- so the
platform's overhead is not the thing to optimize.

The structural claim survives the parameterization: memory tracks the agents
that are \emph{awake}. Under any of these rows, a billion registered agents at
realistic concurrency is a handful of commodity nodes rather than a datacenter,
and the difference between this and the 128\,TB estimate in \S\ref{sec:intro} is
entirely that the estimate charged every registered agent for an isolation
boundary it was not using.

\subsection{What the bill looks like}

At published rates for the always-on control plane --- \rPlaneCpu{}\,vCPU and
\rPlaneRam{}\,GB, \$\rPlaneCost{}/month~\cite{railway} \cited{} --- the fixed
cost of the platform exceeds the storage cost of a billion dormant agents by two
orders of magnitude. In other words, at any population, \emph{the agents are not
the bill}. The bill is the compute for the agents that are awake, plus the item
in \S\ref{sec:limits} that grows without bound.

We state one exclusion emphatically, because omitting it would make every figure
here misleading: \textbf{model tokens are not in this model and they dominate
it.} A single frontier-model call costs cents; the infrastructure under an agent
turn costs microdollars. Any per-call price derived from this paper sits
\emph{beside} token billing and must never be presented as including it.
